\documentclass[12pt,a4paper]{report}

\usepackage[utf8]{inputenc}
\usepackage[T1]{fontenc}
\usepackage[french]{babel}
\usepackage[table]{xcolor}
\usepackage{longtable}
\usepackage{listings}
\usepackage{hyperref}
\usepackage{array}
\usepackage{geometry}
\usepackage{enumitem}
\usepackage{amssymb}
\usepackage{textcomp}

\hypersetup{
    colorlinks=true,
    linkcolor=blue,
    urlcolor=blue,
}

\geometry{margin=2.5cm}

\lstset{
    basicstyle=\footnotesize\ttfamily,
    breaklines=true,
    frame=single,
    numbers=left,
    numberstyle=\tiny,
    inputencoding=utf8,
    extendedchars=true,
    literate=
        {é}{{\'e}}1
        {è}{{\`e}}1
        {ê}{{\^e}}1
        {ë}{{\"e}}1
        {à}{{\`a}}1
        {â}{{\^a}}1
        {ä}{{\"a}}1
        {ù}{{\`u}}1
        {û}{{\^u}}1
        {ü}{{\"u}}1
        {ô}{{\^o}}1
        {ö}{{\"o}}1
        {î}{{\^i}}1
        {ï}{{\"i}}1
        {ç}{{\c c}}1
        {É}{{\'E}}1
        {È}{{\`E}}1
        {Ê}{{\^E}}1
        {À}{{\`A}}1
        {Ç}{{\c C}}1
        {€}{{\texteuro}}1
        {—}{{---}}1
        {→}{{$\rightarrow$}}1
        {←}{{$\leftarrow$}}1
        {─}{{-}}1
        {┌}{{+}}1
        {├}{{+}}1
        {┼}{{+}}1
        {└}{{+}}1
}

\begin{document}

\begin{titlepage}
\begin{center}
\vspace*{5cm}
{\Huge \textbf{Travaux Pratiques}}\\[1cm]
{\LARGE SEO Expert}\\[2cm]
\rule{\textwidth}{0.5pt}\\[1cm]
{\large Document regroupant l'ensemble des travaux pratiques}\\
{\large pour la formation SEO Expert}\\[2cm]
\rule{\textwidth}{0.5pt}\\[2cm]
{\large \today}
\end{center}
\end{titlepage}

\tableofcontents
\newpage

%%%%%%%%%%%%%%%%%%%%%%%%%%%%%%%%%%%%%%%%%%%%%%%%%%%
\chapter{TP 01 --- Audit SEO Complet d'un Site Web}
%%%%%%%%%%%%%%%%%%%%%%%%%%%%%%%%%%%%%%%%%%%%%%%%%%%

\section{Objectifs}
\begin{itemize}
\item Réaliser un audit SEO technique complet
\item Utiliser les outils professionnels (Screaming Frog, Google Search Console, Lighthouse)
\item Produire un rapport d'audit avec préconisations actionnables
\end{itemize}

\section{Outils nécessaires}
\begin{itemize}
\item \href{https://www.screamingfrog.co.uk/seo-spider/}{Screaming Frog SEO Spider} (version gratuite)
\item Google Search Console (ou rapport de démonstration fourni)
\item Google PageSpeed Insights / Lighthouse
\item Mobile-Friendly Test
\item Extension Web Developer Toolbar
\end{itemize}

\section{Exercice 1 --- Crawl du site avec Screaming Frog (40 min)}

\textbf{Consigne :} Lancez Screaming Frog sur un site de votre choix (votre propre site ou un site fictif fourni par le formateur).

\textbf{Étapes :}
\begin{enumerate}
\item Configurer le crawl : agents user-agent ``Googlebot'', respecter robots.txt
\item Lancer le crawl (limitez à 500 URLs pour la version gratuite)
\item Exporter les résultats
\end{enumerate}

\textbf{À analyser :}

\begin{tabular}{|l|c|l|}
\hline
\textbf{Métrique} & \textbf{Valeur} & \textbf{Diagnostic} \\
\hline
Nombre total d'URLs crawlées & & \\
\hline
URLs OK (200) & & \\
\hline
URLs 301 (redirections) & & \\
\hline
URLs 404 (non trouvées) & & \\
\hline
URLs 5XX (erreurs serveur) & & \\
\hline
URLs avec meta robots noindex & & \\
\hline
URLs avec canonical & & \\
\hline
URLs avec title dupliqués & & \\
\hline
URLs sans title & & \\
\hline
URLs sans meta description & & \\
\hline
Temps de réponse moyen & & \\
\hline
\end{tabular}

\textbf{Livrable :} Tableau récapitulatif avec les actions correctives pour chaque problème.

\section{Exercice 2 --- Analyse des balises On-Page (30 min)}

\textbf{Consigne :} Utilisez les exports de Screaming Frog pour analyser la qualité des balises.

\textbf{1. Analyse des Title Tags :}
\begin{itemize}
\item Combien de pages ont un title $>$ 60 caractères ?
\item Combien de pages ont un title $<$ 30 caractères ?
\item Combien de titles sont dupliqués ?
\item Proposez des corrections pour 5 titles problématiques.
\end{itemize}

\textbf{2. Analyse des Meta Descriptions :}
\begin{itemize}
\item Combien de pages ont une meta description $>$ 160 caractères ?
\item Combien de pages n'ont pas de meta description ?
\item Proposez 5 meta descriptions optimisées.
\end{itemize}

\textbf{3. Analyse des Headings :}
\begin{itemize}
\item Combien de pages ont plus d'un H1 ?
\item Combien de pages n'ont pas de H1 ?
\item La structure H1 $\rightarrow$ H2 $\rightarrow$ H3 est-elle logique sur les pages principales ?
\end{itemize}

\section{Exercice 3 --- Analyse des Core Web Vitals (30 min)}

\textbf{Consigne :} Testez 5 pages du site avec Google PageSpeed Insights.

\begin{tabular}{|l|c|c|c|c|c|}
\hline
\textbf{Page} & \textbf{LCP} & \textbf{INP} & \textbf{CLS} & \textbf{Score Mobile} & \textbf{Score Desktop} \\
\hline
Accueil & & & & & \\
\hline
Catégorie & & & & & \\
\hline
Article & & & & & \\
\hline
Contact & & & & & \\
\hline
Produit & & & & & \\
\hline
\end{tabular}

\textbf{Rendez-vous :}
\begin{enumerate}
\item Quelles sont les pages les plus lentes ? Pourquoi ?
\item Quelles sont les recommandations de PageSpeed Insights ?
\item Classez les corrections par priorité (rapides vs complexes).
\item Quel est l'impact potentiel sur le classement SEO ?
\end{enumerate}

\section{Exercice 4 --- Analyse du maillage interne (25 min)}

\textbf{Consigne :} Dans Screaming Frog, allez dans l'onglet ``Internal Links''.

\begin{enumerate}
\item Combien de pages ont moins de 3 liens internes entrants ?
\item Identifiez les pages orphelines (0 liens internes entrants).
\item Quelle est la profondeur de clics moyenne du site ?
\item Quelles pages sont les plus liées (hubs) ?
\item Proposez un plan de maillage interne pour corriger les problèmes.
\end{enumerate}

\textbf{Visualisation :} Utilisez la vue ``Force-Directed Graph'' de Screaming Frog (si disponible) pour visualiser le maillage.

\section{Exercice 5 --- Vérification de l'indexation (20 min)}

\textbf{Consigne :} Vérifiez l'indexation du site avec 3 méthodes différentes.

\textbf{Méthode 1 --- Google Search Console :}
\begin{itemize}
\item Combien de pages sont indexées ?
\item Y a-t-il des erreurs de couverture ?
\item Quelles pages sont exclues et pourquoi ?
\end{itemize}

\textbf{Méthode 2 --- Commande site:}
\begin{lstlisting}
site:exemple.com
\end{lstlisting}
\begin{itemize}
\item Combien de résultats Google retourne-t-il ?
\item Comparez avec le nombre de pages crawlées par Screaming Frog.
\end{itemize}

\textbf{Méthode 3 --- URL Inspection Tool :}
\begin{itemize}
\item Testez 5 URLs spécifiques.
\item Sont-elles indexées ? Si non, pourquoi ?
\end{itemize}

\section{Exercice 6 --- Analyse des concurrents (25 min)}

\textbf{Consigne :} Analysez 3 concurrents directs et comparez avec le site audité.

\begin{tabular}{|l|c|c|c|c|}
\hline
\textbf{Critère} & \textbf{Site audité} & \textbf{Concurrent 1} & \textbf{Concurrent 2} & \textbf{Concurrent 3} \\
\hline
Pages indexées (site:) & & & & \\
\hline
Vitesse mobile & & & & \\
\hline
Balises title optimisées & & & & \\
\hline
Rich snippets & & & & \\
\hline
Backlinks (domaines) & & & & \\
\hline
Trafic estimé & & & & \\
\hline
Mots-clés positionnés & & & & \\
\hline
\end{tabular}

\textbf{Outils recommandés :} Ubersuggest (gratuit), SimilarWeb, Wappalyzer.

\section{Exercice 7 --- Rédaction du rapport d'audit (30 min)}

\textbf{Consigne :} Rédigez un rapport d'audit SEO complet structuré comme suit :

\begin{enumerate}
\item \textbf{Résumé exécutif} (1 page max)
   \begin{itemize}
   \item Score global du site (A/B/C/D/E)
   \item 3 problèmes critiques + 3 recommandations prioritaires
   \item Estimation du gain de trafic potentiel
   \end{itemize}

\item \textbf{Audit technique} (Screaming Frog, GSC)
   \begin{itemize}
   \item Analyse du crawl et de l'indexation
   \item Erreurs 4xx/5xx
   \item Vitesse et Core Web Vitals
   \item Mobile-friendliness
   \end{itemize}

\item \textbf{Audit On-Page}
   \begin{itemize}
   \item Analyse des balises
   \item Qualité du contenu
   \item Maillage interne
   \end{itemize}

\item \textbf{Audit Off-Page}
   \begin{itemize}
   \item Profil de backlinks
   \item Présence dans les annuaires
   \item E-E-A-T
   \end{itemize}

\item \textbf{Plan d'action}
   \begin{itemize}
   \item Priorité haute (urgent, corrigé sous 1 semaine)
   \item Priorité moyenne (corrigé sous 1 mois)
   \item Priorité basse (à planifier)
   \item Estimation du temps et des ressources nécessaires
   \end{itemize}
\end{enumerate}

\section{Exercice 8 --- Présentation orale (15 min par groupe)}

\textbf{Consigne :} Présentez votre audit à la classe en 15 min :
\begin{enumerate}
\item Problèmes principaux (3 min)
\item Recommandations clés (5 min)
\item Plan d'action priorisé (5 min)
\item Questions (2 min)
\end{enumerate}

\textbf{Critères d'évaluation :}
\begin{itemize}
\item Pertinence des problèmes identifiés
\item Qualité des recommandations
\item Priorisation des actions
\item Clarté de la présentation
\item Réponses aux questions
\end{itemize}

\section{Rendu final}

\textbf{À rendre :}
\begin{enumerate}
\item Rapport d'audit PDF (5-10 pages)
\item Présentation (5-10 slides)
\item Données brutes (export Screaming Frog + captures d'écran)
\end{enumerate}

\textbf{Barème :}
\begin{itemize}
\item Rapport écrit : /20
\item Présentation : /10
\item Pertinence des préconisations : /10
\item Total : /40
\end{itemize}

%%%%%%%%%%%%%%%%%%%%%%%%%%%%%%%%%%%%%%%%%%%%%%%%%%%
\chapter{TP 02 --- Stratégie de Contenu et Keyword Research}
%%%%%%%%%%%%%%%%%%%%%%%%%%%%%%%%%%%%%%%%%%%%%%%%%%%

\section{Objectifs}
\begin{itemize}
\item Réaliser une recherche de mots-clés complète avec des outils professionnels
\item Créer un plan de contenu stratégique (topical clusters)
\item Rédiger et optimiser un article de blog SEO
\end{itemize}

\section{Outils nécessaires}
\begin{itemize}
\item Google Keyword Planner (compte Google Ads requis)
\item Ubersuggest (version gratuite)
\item AnswerThePublic
\item Google Trends
\item SEMrush / Ahrefs (si disponible)
\item Google Docs ou un CMS (WordPress, etc.)
\end{itemize}

\section{Exercice 1 --- Définition du persona et du parcours client (20 min)}

\textbf{Consigne :} Définissez le client idéal pour votre site.

\textbf{Persona :}
\begin{itemize}
\item Nom fictif : \_\_\_\_\_
\item Âge : \_\_\_\_\_
\item Profession : \_\_\_\_\_
\item Niveau de connaissance en SEO : Débutant / Intermédiaire / Avancé
\item Objectifs : \_\_\_\_\_
\item Points de douleur : \_\_\_\_\_
\item Supports de recherche favoris : \_\_\_\_\_
\end{itemize}

\textbf{Parcours client (User Journey) :}

\begin{tabular}{|l|l|l|l|l|}
\hline
\textbf{Étape} & \textbf{Intention} & \textbf{Question posée} & \textbf{Type de contenu} & \textbf{Canal} \\
\hline
Découverte & Informationnelle & & & \\
\hline
Considération & Commerciale & & & \\
\hline
Comparaison & Commerciale & & & \\
\hline
Décision & Transactionnelle & & & \\
\hline
Fidélisation & Informationnelle & & & \\
\hline
\end{tabular}

\section{Exercice 2 --- Recherche de mots-clés : 200+ mots-clés (45 min)}

\textbf{Consigne :} Générez une liste de 200 mots-clés minimum en utilisant les techniques suivantes.

\textbf{Technique 1 --- Brainstorming (20 mots) :}
\begin{lstlisting}
[Listez 20 mots-cles de base en 5 min]
\end{lstlisting}

\textbf{Technique 2 --- Google Suggest (30 mots) :}
\begin{itemize}
\item Tapez vos mots-clés de base dans Google
\item Notez les suggestions automatiques
\item Utilisez les lettres A-Z après chaque mot-clé
\end{itemize}

\textbf{Technique 3 --- ``Questions'' (20 mots) :}
\begin{itemize}
\item Utilisez AnswerThePublic
\item Notez les questions, prépositions, comparaisons
\item Catégorisez par intention
\end{itemize}

\textbf{Technique 4 --- Google Keyword Planner (50+ mots) :}
\begin{itemize}
\item Utilisez le Keyword Planner
\item Entrez vos 10 meilleurs mots-clés de base
\item Exportez les suggestions + données de volume
\end{itemize}

\textbf{Technique 5 --- Ubersuggest / SEMrush (50+ mots) :}
\begin{itemize}
\item Entrez un domaine concurrent
\item Récupérez les mots-clés de vos concurrents
\item Identifiez les lacunes (gaps)
\end{itemize}

\textbf{Technique 6 --- Google Trends (30 mots) :}
\begin{itemize}
\item Comparez la saisonnalité des mots-clés
\item Identifiez les tendances émergentes
\item Notez les requêtes en hausse
\end{itemize}

\textbf{Livrable :} Un fichier CSV/Google Sheets avec :

\begin{tabular}{|l|c|c|c|l|c|}
\hline
\textbf{Mot-clé} & \textbf{Volume mensuel} & \textbf{Difficulté} & \textbf{Intention} & \textbf{Page cible} & \textbf{Priorité} \\
\hline
 & & & & & \\
\hline
\end{tabular}

\section{Exercice 3 --- Analyse de l'intention de recherche (25 min)}

\textbf{Consigne :} Pour 10 mots-clés de votre liste, analysez l'intention via la SERP.

\textbf{Méthode :} Tapez chaque mot-clé dans Google (navigation privée) et analysez les résultats.

\begin{longtable}{|l|l|l|l|}
\hline
\textbf{Mot-clé} & \textbf{Type de résultats dominants (articles, produits, vidéos...)} & \textbf{Intention réelle} & \textbf{Est-ce que l'intention supposée était correcte ?} \\
\hline
1 & & & \\
\hline
2 & & & \\
\hline
3 & & & \\
\hline
4 & & & \\
\hline
5 & & & \\
\hline
6 & & & \\
\hline
7 & & & \\
\hline
8 & & & \\
\hline
9 & & & \\
\hline
10 & & & \\
\hline
\end{longtable}

\textbf{Questions :}
\begin{enumerate}
\item Quels mots-clés ont une intention mixte ?
\item Comment adapter votre contenu pour les intentions mixtes ?
\item Y a-t-il des featured snippets sur vos mots-clés ? Pour quel type de contenu ?
\end{enumerate}

\section{Exercice 4 --- Création d'un Topical Cluster (30 min)}

\textbf{Consigne :} Créez une structure de Topical Cluster basée sur vos mots-clés.

\textbf{Page pilier (pillar page) :}
\begin{itemize}
\item Sujet principal : \_\_\_\_\_
\item URL cible : \_\_\_\_\_
\item Mots-clés principaux : \_\_\_\_\_
\item Structure H2/H3 de la page pilier :
   \begin{itemize}
   \item H2 : \_\_\_\_\_
      \begin{itemize}
      \item H3 : \_\_\_\_\_
      \item H3 : \_\_\_\_\_
      \end{itemize}
   \item H2 : \_\_\_\_\_
      \begin{itemize}
      \item H3 : \_\_\_\_\_
      \end{itemize}
   \item H2 : \_\_\_\_\_
   \end{itemize}
\end{itemize}

\textbf{Articles du cluster (10 articles minimum) :}

\begin{longtable}{|l|l|c|c|l|l|}
\hline
\textbf{Article} & \textbf{Mot-clé cible} & \textbf{Volume} & \textbf{Intention} & \textbf{Lien vers pilier} & \textbf{Lien vers article suivant} \\
\hline
1 & & & & & \\
\hline
2 & & & & & \\
\hline
3 & & & & & \\
\hline
4 & & & & & \\
\hline
5 & & & & & \\
\hline
6 & & & & & \\
\hline
7 & & & & & \\
\hline
8 & & & & & \\
\hline
9 & & & & & \\
\hline
10 & & & & & \\
\hline
\end{longtable}

\textbf{Schéma des liens :}
\begin{lstlisting}[basicstyle=\ttfamily\small]
                    +-- Article 1
                    +-- Article 2
  Page Pilier ------+-- Article 3  <-> Article 4
                    +-- Article 5
                    +-- Article 6  <-> Article 7
                    +-- ...
                    +-- Article 10
\end{lstlisting}

\section{Exercice 5 --- Rédaction et Optimisation d'un Article SEO (50 min)}

\textbf{Consigne :} Rédigez un article de blog complet (1500-2000 mots) optimisé SEO.

\textbf{Structure imposée :}
\begin{itemize}
\item Title (50-60 caractères) : \_\_\_\_\_
\item Meta description (150-160 caractères) : \_\_\_\_\_
\item URL slug (3-5 mots) : \_\_\_\_\_
\item H1 (contenant le mot-clé principal) : \_\_\_\_\_
\item Introduction (100-150 mots, mot-clé dans les 100 premiers mots)
\item 4-6 sections H2 (chacune avec mot-clé secondaire)
\item Sous-sections H3 si nécessaire
\item Conclusion (50-100 mots) avec appel à l'action
\item 3-5 liens internes vers d'autres articles du cluster
\item 2-3 liens externes vers des sources faisant autorité
\item 1-2 images avec alt text optimisé
\end{itemize}

\textbf{Checklist d'optimisation :}

\begin{tabular}{|l|c|l|}
\hline
\textbf{Critère} & \textbf{OK} & \textbf{Remarque} \\
\hline
Mot-clé principal dans les 100 premiers mots & & \\
\hline
Mot-clé principal dans l'URL & & \\
\hline
Mot-clé principal dans le H1 & & \\
\hline
Mots-clés secondaires dans les H2 & & \\
\hline
Densité de mot-clé naturelle (pas de bourrage) & & \\
\hline
Balise title optimisée & & \\
\hline
Meta description optimisée & & \\
\hline
Images avec alt text & & \\
\hline
Liens internes ($\ge$3) & & \\
\hline
Liens externes ($\ge$2) & & \\
\hline
Article $\ge$ 1500 mots & & \\
\hline
Appel à l'action présent & & \\
\hline
Données structurées (Article) & & \\
\hline
Temps de lecture $<$ 15 min & & \\
\hline
\end{tabular}

\section{Exercice 6 --- Analyse de la concurrence SERP (25 min)}

\textbf{Consigne :} Analysez le top 5 des résultats Google pour votre mot-clé principal.

\begin{longtable}{|c|l|l|c|c|c|c|l|l|}
\hline
\textbf{Position} & \textbf{URL} & \textbf{Title} & \textbf{Longueur} & \textbf{Mots-clés} & \textbf{Backlinks (domaines)} & \textbf{Âge domaine} & \textbf{Points forts} & \textbf{Points faibles} \\
\hline
1 & & & & & & & & \\
\hline
2 & & & & & & & & \\
\hline
3 & & & & & & & & \\
\hline
4 & & & & & & & & \\
\hline
5 & & & & & & & & \\
\hline
\end{longtable}

\textbf{Analyse des lacunes (Content Gap) :}
\begin{enumerate}
\item Quels sujets les concurrents ne couvrent-ils pas ?
\item Quels angles originaux pouvez-vous apporter ?
\item Quel format (vidéo, infographie, étude de cas) manque dans la SERP ?
\item Comment votre contenu sera-t-il meilleur que le \#1 actuel ?
\end{enumerate}

\section{Exercice 7 --- Création d'un Calendrier Éditorial (20 min)}

\textbf{Consigne :} Planifiez 3 mois de contenu SEO.

\begin{longtable}{|l|l|l|l|l|l|}
\hline
\textbf{Semaine} & \textbf{Date} & \textbf{Article} & \textbf{Mot-clé cible} & \textbf{Auteur} & \textbf{Statut} \\
\hline
S1 & & & & & \\
\hline
S2 & & & & & \\
\hline
S3 & & & & & \\
\hline
S4 & & & & & \\
\hline
S5 & & & & & \\
\hline
\ldots & & & & & \\
\hline
\end{longtable}

\textbf{Objectifs du calendrier :}
\begin{itemize}
\item 2-3 articles par semaine minimum
\item 1 page pilier par mois
\item 1 mise à jour de contenu existant par mois
\item 1 contenu ``Big Rock'' (étude de cas, guide ultime) par trimestre
\end{itemize}

\section{Exercice 8 --- Promotion de contenu (15 min)}

\textbf{Consigne :} Planifiez la promotion de votre article.

\begin{tabular}{|l|l|l|l|}
\hline
\textbf{Canal} & \textbf{Action} & \textbf{Date} & \textbf{KPI} \\
\hline
Email (newsletter) & Envoyer aux abonnés & & Taux d'ouverture, clics \\
\hline
LinkedIn & Publier un post & & Impressions, engagement \\
\hline
Twitter/X & Publier un thread & & Vues, clics \\
\hline
Facebook & Partager dans les groupes & & Reach, clics \\
\hline
Pinterest & Épingler avec image & & Impressions \\
\hline
Forums (Reddit, Quora) & Répondre avec lien & & Upvotes, trafic \\
\hline
Backlinks & Outreach aux sites partenaires & & Liens obtenus \\
\hline
Webinar / Live & Présenter le sujet & & Participants \\
\hline
\end{tabular}

\section{Rendu final}

\textbf{À rendre :}
\begin{enumerate}
\item Fichier CSV des mots-clés (200+ entrées)
\item Structure du topical cluster (schéma + pages)
\item Article de blog complet optimisé (Google Doc ou HTML)
\item Calendrier éditorial 3 mois
\item Plan de promotion
\end{enumerate}

\textbf{Barème :}
\begin{itemize}
\item Recherche de mots-clés : /10
\item Structure du cluster : /10
\item Qualité de l'article : /20
\item Calendrier + promotion : /10
\item Total : /50
\end{itemize}

%%%%%%%%%%%%%%%%%%%%%%%%%%%%%%%%%%%%%%%%%%%%%%%%%%%
\chapter{TP 03 --- Optimisation Technique et Performance Web}
%%%%%%%%%%%%%%%%%%%%%%%%%%%%%%%%%%%%%%%%%%%%%%%%%%%

\section{Objectifs}
\begin{itemize}
\item Optimiser la vitesse et les Core Web Vitals d'un site web
\item Mettre en place les données structurées (Schema Markup)
\item Configurer robots.txt, sitemap XML et balises techniques
\item Implémenter une stratégie de rendu optimisée pour le SEO
\end{itemize}

\section{Outils nécessaires}
\begin{itemize}
\item Serveur local (XAMPP, WAMP, ou environnement en ligne)
\item Éditeur de code (VS Code)
\item Google PageSpeed Insights / Lighthouse
\item Schema Markup Validator (Google)
\item Screaming Frog (ou équivalent)
\end{itemize}

\section{Exercice 1 --- Analyse de Performance Initiale (20 min)}

\textbf{Consigne :} Avant toute optimisation, mesurez l'état initial du site.

\textbf{PageSpeed Insights :}
\begin{tabular}{|l|c|c|c|}
\hline
\textbf{Métrique} & \textbf{Valeur initiale} & \textbf{Objectif} & \textbf{Statut} \\
\hline
LCP & & $\le$ 2.5s & \\
\hline
INP & & $\le$ 200ms & \\
\hline
CLS & & $\le$ 0.1 & \\
\hline
FCP & & $\le$ 1.8s & \\
\hline
TTFB & & $\le$ 0.8s & \\
\hline
Score Mobile & & $\ge$ 90 & \\
\hline
Score Desktop & & $\ge$ 90 & \\
\hline
\end{tabular}

\textbf{Lighthouse (audit complet) :}
\begin{tabular}{|l|c|l|}
\hline
\textbf{Catégorie} & \textbf{Score} & \textbf{Problèmes principaux} \\
\hline
Performance & & \\
\hline
Accessibilité & & \\
\hline
Best Practices & & \\
\hline
SEO & & \\
\hline
PWA & & \\
\hline
\end{tabular}

\textbf{Capture d'écran avant optimisation :} Prenez une capture de l'écran PageSpeed Insights.

\section{Exercice 2 --- Optimisation des Images (30 min)}

\textbf{Consigne :} Optimisez toutes les images du site.

\textbf{Inventaire des images :}
\begin{tabular}{|l|l|c|c|l|l|}
\hline
\textbf{Image} & \textbf{Format} & \textbf{Taille (Ko)} & \textbf{Dimensions} & \textbf{Alt text} & \textbf{Action} \\
\hline
 & & & & & \\
\hline
 & & & & & \\
\hline
\end{tabular}

\textbf{Actions à réaliser :}
\begin{enumerate}
\item \textbf{Conversion en format moderne :} PNG/JPG $\rightarrow$ WebP ou AVIF
\item \textbf{Compression :} Utilisez Squoosh, TinyPNG ou ImageOptim
\item \textbf{Redimensionnement :} Ajustez aux dimensions d'affichage réelles
\item \textbf{Lazy loading :} Ajoutez \texttt{loading="lazy"} sur les images hors écran
\item \textbf{Responsive images :} Utilisez \texttt{srcset} et \texttt{sizes} pour les écrans Retina
\item \textbf{Alt text :} Ajoutez des descriptions pertinentes pour chaque image
\end{enumerate}

\textbf{Code à générer :}
\begin{lstlisting}[language=HTML]
<!-- Exemple d'image responsive optimisee -->
<img 
  src="image-800w.webp"
  srcset="image-400w.webp 400w, image-800w.webp 800w, image-1200w.webp 1200w"
  sizes="(max-width: 600px) 400px, (max-width: 1200px) 800px, 1200px"
  loading="lazy"
  alt="Description SEO de l'image"
  width="800"
  height="600"
  decoding="async"
>
\end{lstlisting}

\section{Exercice 3 --- Optimisation du CSS et JavaScript (35 min)}

\textbf{Consigne :} Optimisez le rendu du CSS et du JS.

\textbf{1. Critical CSS :}
Identifiez le CSS nécessaire au rendu initial (``above the fold'') :
\begin{lstlisting}[language=HTML]
<!-- Critical CSS en inline dans le <head> -->
<style>
  /* Styles critiques pour le rendu initial */
  body { font-family: ... }
  header { ... }
  .hero { ... }
</style>
<!-- Chargement differe du CSS complet -->
<link rel="preload" href="styles.css" as="style" onload="this.onload=null;this.rel='stylesheet'">
<noscript><link rel="stylesheet" href="styles.css"></noscript>
\end{lstlisting}

\textbf{2. Minification :}
\begin{itemize}
\item CSS : Utilisez Clean-CSS ou cssnano
\item JS : Utilisez Terser ou UglifyJS
\item HTML : Utilisez html-minifier
\end{itemize}

\textbf{3. Code Splitting :}
\begin{lstlisting}[language=JavaScript]
// Exemple : chargement dynamique des modules JS
import('heavy-module.js').then(module => {
  module.init();
});
\end{lstlisting}

\textbf{4. Suppression du render-blocking :}
\begin{itemize}
\item Déplacez les scripts non critiques en bas du \texttt{<body>}
\item Ajoutez \texttt{async} ou \texttt{defer} sur les scripts
\end{itemize}
\begin{lstlisting}[language=HTML]
<script src="analytics.js" async></script>
<script src="app.js" defer></script>
\end{lstlisting}

\section{Exercice 4 --- Optimisation du Serveur (25 min)}

\textbf{Consigne :} Configurez le serveur pour de meilleures performances.

\textbf{1. Cache navigateur (fichier .htaccess pour Apache ou configuration Nginx) :}
\begin{lstlisting}
# Apache .htaccess
<IfModule mod_expires.c>
  ExpiresActive On
  ExpiresByType image/webp "access plus 1 year"
  ExpiresByType image/avif "access plus 1 year"
  ExpiresByType text/css "access plus 1 month"
  ExpiresByType application/javascript "access plus 1 month"
  ExpiresByType text/html "access plus 0 seconds"
</IfModule>
\end{lstlisting}

\textbf{2. Compression Brotli / Gzip :}
\begin{lstlisting}
# Apache
AddOutputFilterByType BROTLI_COMPRESS text/html text/css application/javascript
AddOutputFilterByType DEFLATE text/html text/css application/javascript
\end{lstlisting}

\textbf{3. Headers de cache et sécurité :}
\begin{lstlisting}
<IfModule mod_headers.c>
  Header set X-Content-Type-Options "nosniff"
  Header set X-Frame-Options "SAMEORIGIN"
  Header set Referrer-Policy "strict-origin-when-cross-origin"
  # Cache-Control pour les ressources statiques
  Header set Cache-Control "public, max-age=31536000, immutable" "exprès:image/*"
</IfModule>
\end{lstlisting}

\textbf{Questions :}
\begin{enumerate}
\item Quelle est la différence entre Gzip et Brotli ?
\item Pourquoi les ressources statiques doivent-elles avoir un cache long ?
\item Qu'est-ce que l'ETag et à quoi sert-il ?
\end{enumerate}

\section{Exercice 5 --- Mise en place des Données Structurées (35 min)}

\textbf{Consigne :} Implémentez les données structurées JSON-LD sur le site.

\textbf{1. Organisation :}
\begin{lstlisting}[language=JSON]
{
  "@context": "https://schema.org",
  "@type": "Organization",
  "name": "Nom de l'entreprise",
  "url": "https://exemple.com",
  "logo": "https://exemple.com/logo.png",
  "contactPoint": {
    "@type": "ContactPoint",
    "telephone": "+33-1-23-45-67-89",
    "contactType": "customer service"
  },
  "sameAs": [
    "https://www.facebook.com/entreprise",
    "https://www.linkedin.com/company/entreprise"
  ]
}
\end{lstlisting}

\textbf{2. Article / BlogPost :}
Implémentez les données structurées pour un article :
\begin{lstlisting}[language=JSON]
{
  "@context": "https://schema.org",
  "@type": "Article",
  "headline": "Titre de l'article",
  "description": "Description de l'article",
  "author": {
    "@type": "Person",
    "name": "Nom de l'auteur"
  },
  "datePublished": "2026-01-15",
  "dateModified": "2026-01-20",
  "image": "https://exemple.com/image.jpg",
  "publisher": {
    "@type": "Organization",
    "name": "Nom du site",
    "logo": {
      "@type": "ImageObject",
      "url": "https://exemple.com/logo.png"
    }
  }
}
\end{lstlisting}

\textbf{3. BreadcrumbList :}
\begin{lstlisting}[language=JSON]
{
  "@context": "https://schema.org",
  "@type": "BreadcrumbList",
  "itemListElement": [
    { "@type": "ListItem", "position": 1, "name": "Accueil", "item": "https://exemple.com" },
    { "@type": "ListItem", "position": 2, "name": "Blog", "item": "https://exemple.com/blog" },
    { "@type": "ListItem", "position": 3, "name": "Article SEO", "item": "https://exemple.com/blog/article-seo" }
  ]
}
\end{lstlisting}

\textbf{4. FAQPage :}
\begin{lstlisting}[language=JSON]
{
  "@context": "https://schema.org",
  "@type": "FAQPage",
  "mainEntity": [
    {
      "@type": "Question",
      "name": "Qu'est-ce que le SEO ?",
      "acceptedAnswer": {
        "@type": "Answer",
        "text": "Le SEO (Search Engine Optimization) est l'ensemble des techniques..."
      }
    }
  ]
}
\end{lstlisting}

\textbf{Validation :} Testez chaque implémentation avec le \href{https://validator.schema.org/}{Schema Markup Validator} de Google.

\section{Exercice 6 --- Configuration Technique du Site (30 min)}

\textbf{Consigne :} Mettez en place les fichiers techniques essentiels.

\textbf{1. Robots.txt :}
\begin{lstlisting}
User-agent: *
Disallow: /admin/
Disallow: /private/
Allow: /admin/public/

Sitemap: https://exemple.com/sitemap.xml
\end{lstlisting}

\textbf{À faire :}
\begin{itemize}
\item Testez le robots.txt via Google Search Console
\item Vérifiez que les pages importantes ne sont pas bloquées
\end{itemize}

\textbf{2. Sitemap XML :}
\begin{lstlisting}[language=XML]
<?xml version="1.0" encoding="UTF-8"?>
<urlset xmlns="http://www.sitemaps.org/schemas/sitemap/0.9">
  <url>
    <loc>https://exemple.com/</loc>
    <lastmod>2026-01-15</lastmod>
    <changefreq>weekly</changefreq>
    <priority>1.0</priority>
  </url>
  <!-- Repeter pour chaque page -->
</urlset>
\end{lstlisting}

\textbf{À faire :}
\begin{itemize}
\item Générez le sitemap dynamiquement ou avec un outil
\item Soumettez-le via Google Search Console
\item Vérifiez qu'il ne contient pas d'URLs bloquées par robots.txt
\end{itemize}

\textbf{3. Fichier .htaccess (redirections 301) :}
\begin{lstlisting}
# Redirection des URLs www vers non-www (ou inverse)
RewriteCond %{HTTP_HOST} ^www\.exemple\.com [NC]
RewriteRule ^(.*)$ https://exemple.com/$1 [L,R=301]

# Redirection HTTP vers HTTPS
RewriteCond %{HTTPS} off
RewriteRule ^(.*)$ https://%{HTTP_HOST}/$1 [L,R=301]
\end{lstlisting}

\section{Exercice 7 --- Optimisation Mobile (20 min)}

\textbf{Consigne :} Assurez la compatibilité mobile du site.

\textbf{Checklist Mobile-First :}
\begin{itemize}
\item[$\square$] Viewport meta tag présent : \texttt{<meta name="viewport" content="width=device-width, initial-scale=1">}
\item[$\square$] Police de caractères $\ge$ 16px pour le texte
\item[$\square$] Éléments cliquables $\ge$ 48$\times$48px
\item[$\square$] Pas de contenu horizontal (scroll horizontal)
\item[$\square$] Images responsives (srcset)
\item[$\square$] Pas de popups intrusifs (interstitiels)
\item[$\square$] Test Mobile-Friendly Google réussi
\item[$\square$] Temps de chargement mobile $<$ 3s
\end{itemize}

\textbf{Test :}
\begin{enumerate}
\item Utilisez l'outil Mobile-Friendly Test de Google
\item Testez la navigation mobile sur 5 pages
\item Capturez les problèmes d'affichage
\item Corrigez les problèmes identifiés
\end{enumerate}

\section{Exercice 8 --- Test Final et Comparaison (20 min)}

\textbf{Consigne :} Après toutes les optimisations, re-mesurez les performances.

\textbf{Comparatif avant/après :}

\begin{tabular}{|l|c|c|c|}
\hline
\textbf{Métrique} & \textbf{Avant} & \textbf{Après} & \textbf{Amélioration} \\
\hline
LCP & & & \\
\hline
INP & & & \\
\hline
CLS & & & \\
\hline
Score Mobile & & & \\
\hline
Score Desktop & & & \\
\hline
Taille page (Ko) & & & \\
\hline
Requêtes & & & \\
\hline
Temps chargement & & & \\
\hline
\end{tabular}

\textbf{Rapport final :}
\begin{enumerate}
\item Quelles optimisations ont eu le plus d'impact ?
\item Quel est le gain estimé en termes d'expérience utilisateur ?
\item Quel serait l'impact SEO potentiel ?
\item Quelles optimisations reste-t-il à faire (prochain sprint) ?
\end{enumerate}

\textbf{Livrable :} Fiche récapitulative des optimisations réalisées avec captures d'écran avant/après.

\section{Rendu final}

\textbf{À rendre :}
\begin{enumerate}
\item Rapport d'optimisation technique (avant/après avec captures d'écran)
\item Fichiers de configuration produits (robots.txt, .htaccess, sitemap.xml)
\item Code des données structurées JSON-LD
\item Résultat PageSpeed Insights avant vs après
\end{enumerate}

\textbf{Barème :}
\begin{itemize}
\item Analyse initiale : /5
\item Optimisation images et ressources : /10
\item Optimisation serveur et cache : /10
\item Données structurées : /10
\item Configuration technique : /10
\item Résultat final (amélioration mesurable) : /5
\item Total : /50
\end{itemize}

%%%%%%%%%%%%%%%%%%%%%%%%%%%%%%%%%%%%%%%%%%%%%%%%%%%
\chapter{TP 04 --- Projet Final : Campagne SEO Complète}
%%%%%%%%%%%%%%%%%%%%%%%%%%%%%%%%%%%%%%%%%%%%%%%%%%%

\section{Objectifs}
\begin{itemize}
\item Mettre en œuvre une stratégie SEO complète de A à Z
\item Appliquer l'ensemble des compétences acquises (technique, contenu, netlinking)
\item Mesurer et optimiser les résultats sur la durée
\end{itemize}

\section{Présentation du Projet}

Vous êtes consultant SEO en agence. Un nouveau client vous confie la mission d'optimiser son site web pour les moteurs de recherche. Vous avez 4 semaines pour réaliser un audit complet, déployer les optimisations et présenter les résultats.

\textbf{Client fictif (au choix) :}
\begin{itemize}
\item \textbf{Option A :} E-commerce de produits bio (500 produits)
\item \textbf{Option B :} Blog de voyage (200 articles)
\item \textbf{Option C :} Site vitrine d'agence web (20 pages)
\item \textbf{Option D :} Site institutionnel d'association (50 pages)
\end{itemize}

\section{Semaine 1 --- Audit et Diagnostic}

\textbf{Jour 1-2 : Audit technique}
\begin{itemize}
\item[$\square$] Crawl complet avec Screaming Frog
\item[$\square$] Analyse des erreurs 4xx, 5xx, 301
\item[$\square$] Analyse des balises (title, meta, headings)
\item[$\square$] Vérification de l'indexation (GSC + site:)
\item[$\square$] Analyse de la vitesse (PageSpeed Insights)
\item[$\square$] Vérification mobile-friendliness
\end{itemize}

\textbf{Jour 3-4 : Audit du contenu}
\begin{itemize}
\item[$\square$] Inventaire du contenu existant
\item[$\square$] Identification des pages à faible performance
\item[$\square$] Analyse des topical clusters existants
\item[$\square$] Détection de contenu dupliqué
\item[$\square$] Analyse des lacunes (content gaps)
\end{itemize}

\textbf{Jour 5 : Audit Off-Page}
\begin{itemize}
\item[$\square$] Analyse du profil de backlinks
\item[$\square$] Identification des backlinks toxiques
\item[$\square$] Analyse des concurrents (top 5)
\item[$\square$] Évaluation E-E-A-T
\end{itemize}

\textbf{Livrable semaine 1 :}
Rapport d'audit complet avec plan d'action priorisé (impact $\times$ effort).

\section{Semaine 2 --- Optimisations Techniques}

\textbf{Jour 1-2 : Corrections critiques}
\begin{itemize}
\item[$\square$] Correction des erreurs 404
\item[$\square$] Mise en place des redirections 301
\item[$\square$] Correction des boucles de redirection
\item[$\square$] Correction des erreurs 5XX
\end{itemize}

\textbf{Jour 3-4 : Optimisation des performances}
\begin{itemize}
\item[$\square$] Optimisation des images (WebP, compression, lazy loading)
\item[$\square$] Minification CSS/JS
\item[$\square$] Mise en place du cache navigateur
\item[$\square$] Optimisation du TTFB (hébergement, cache serveur)
\end{itemize}

\textbf{Jour 5 : Configuration technique}
\begin{itemize}
\item[$\square$] Robots.txt optimisé
\item[$\square$] Sitemap XML généré et soumis
\item[$\square$] Balises canonical mises en place
\item[$\square$] Redirection www/non-www
\item[$\square$] HTTPS vérifié
\end{itemize}

\textbf{Livrable semaine 2 :}
Fiche des optimisations techniques réalisées + captures d'écran avant/après.

\section{Semaine 3 --- Stratégie de Contenu}

\textbf{Jour 1-2 : Recherche de mots-clés}
\begin{itemize}
\item[$\square$] 200+ mots-clés collectés
\item[$\square$] Analyse de l'intention de recherche
\item[$\square$] Analyse concurrentielle (SERP)
\item[$\square$] Définition des pages piliers et des clusters
\end{itemize}

\textbf{Jour 3-4 : Production de contenu}
\begin{itemize}
\item[$\square$] Rédaction de la page pilier (2000+ mots)
\item[$\square$] Rédaction de 5 articles de cluster (1500+ mots chacun)
\item[$\square$] Optimisation On-Page de chaque article
\item[$\square$] Maillage interne du cluster
\end{itemize}

\textbf{Jour 5 : Optimisation du contenu existant}
\begin{itemize}
\item[$\square$] Mise à jour des 10 pages les plus importantes
\item[$\square$] Ajout de données structurées (Article, FAQ, Organisation)
\item[$\square$] Optimisation des images (alt text)
\item[$\square$] Ajout de liens internes
\end{itemize}

\textbf{Livrable semaine 3 :}
Topical cluster complet + articles publiés + maillage interne.

\section{Semaine 4 --- Netlinking, Suivi et Présentation}

\textbf{Jour 1-2 : Stratégie de netlinking}
\begin{itemize}
\item[$\square$] Identification de 20 sites cibles pour les backlinks
\item[$\square$] Création de contenu ``linkable'' (infographie, étude de cas)
\item[$\square$] Outreach (10 emails envoyés)
\item[$\square$] Guest blogging (1 article invité)
\end{itemize}

\textbf{Jour 3 : Mise en place du suivi}
\begin{itemize}
\item[$\square$] Google Search Console configuré
\item[$\square$] Google Analytics 4 configuré
\item[$\square$] Suivi des positions (rank tracker manuel ou Google Sheets)
\item[$\square$] KPIs définis et dashboard créé
\end{itemize}

\textbf{Jour 4-5 : Rapport final et présentation}

\textbf{Structure du rapport final :}

\begin{enumerate}
\item \textbf{Résumé exécutif} (1 page)
   \begin{itemize}
   \item Situation initiale
   \item Actions menées
   \item Résultats obtenus
   \item Recommandations suite
   \end{itemize}

\item \textbf{Méthodologie} (1 page)
   \begin{itemize}
   \item Outils utilisés
   \item Périmètre de la mission
   \item Calendrier
   \end{itemize}

\item \textbf{Résultats techniques} (2 pages)
   \begin{itemize}
   \item Évolution des Core Web Vitals
   \item Nombre de pages indexées
   \item Erreurs corrigées
   \item Score PageSpeed avant/après
   \end{itemize}

\item \textbf{Résultats contenu} (2 pages)
   \begin{itemize}
   \item Mots-clés ciblés
   \item Contenu créé / optimisé
   \item Évolution du trafic estimé
   \item Topical clusters construits
   \end{itemize}

\item \textbf{Résultats netlinking} (1 page)
   \begin{itemize}
   \item Nouveaux backlinks obtenus
   \item Domaines référents
   \item Évolution du Trust Flow
   \end{itemize}

\item \textbf{KPIs et prochaines étapes} (1 page)
   \begin{itemize}
   \item Tableau de bord des KPIs
   \item Recommandations prochain trimestre
   \item Budget estimé
   \end{itemize}
\end{enumerate}

\section{Exercices complémentaires}

\subsection{Exercice 1 --- Calcul du ROI SEO}
Calculez le retour sur investissement de la campagne :

\begin{lstlisting}
Cout mensuel de la mission : 2 500 EUR
Trafic mensuel avant : 5 000 visites
Trafic mensuel apres : 8 000 visites
Taux de conversion : 2.5%
Panier moyen : 80 EUR
Marge : 40%

1. Quel est le trafic supplementaire ?
2. Combien de conversions supplementaires ?
3. Quel est le chiffre d'affaires supplementaire ?
4. Quelle est la marge supplementaire ?
5. Quel est le ROI ?
\end{lstlisting}

\subsection{Exercice 2 --- Gestion de crise SEO}
Vous recevez une alerte : le site a perdu 50\% de son trafic du jour au lendemain.

\begin{enumerate}
\item Quelles sont les causes possibles ? (listez-en 8 minimum)
\item Quelles vérifications faites-vous en premier ? (ordonnées par priorité)
\item À qui communiquez-vous et que dites-vous ?
\item Quel est votre plan d'action pour les 72 prochaines heures ?
\end{enumerate}

\subsection{Exercice 3 --- Veille concurrentielle}
Configurez un système de veille pour surveiller :

\begin{enumerate}
\item Les nouveaux contenus de vos 3 principaux concurrents
\item Les nouveaux backlinks de vos concurrents
\item Les évolutions de positions des mots-clés
\item Les mises à jour de l'algorithme Google
\item Les nouvelles fonctionnalités SERP (SGE, AI Overviews\ldots)
\end{enumerate}

\textbf{Outils à utiliser :}
\begin{itemize}
\item Google Alerts
\item SEMrush / Ahrefs (alertes)
\item Mention / Talkwalker
\item Feedly / Inoreader (RSS)
\item Twitter / LinkedIn (veille humaine)
\end{itemize}

\section{Grille d'évaluation du projet final}

\begin{tabular}{|l|c|c|c|}
\hline
\textbf{Critère} & \textbf{Poids} & \textbf{Note /10} & \textbf{Commentaire} \\
\hline
Qualité de l'audit initial & 15\% & & \\
\hline
Pertinence des recommandations & 15\% & & \\
\hline
Qualité des optimisations techniques & 15\% & & \\
\hline
Qualité du contenu produit & 15\% & & \\
\hline
Stratégie de netlinking & 10\% & & \\
\hline
Mesure et suivi des résultats & 10\% & & \\
\hline
Rapport final et présentation & 10\% & & \\
\hline
Gestion de projet et respect du planning & 5\% & & \\
\hline
Créativité et initiatives & 5\% & & \\
\hline
\textbf{Total} & \textbf{100\%} & \textbf{/10} & \\
\hline
\end{tabular}

\section{Rendu final}

\textbf{À rendre à la fin de la semaine 4 :}
\begin{enumerate}
\item Rapport de mission complet (10-20 pages)
\item Présentation orale (10-15 slides, 20 min)
\item Fichiers techniques (robots.txt, sitemap, schema, .htaccess)
\item Dashboard de suivi des KPIs (Google Sheets ou Looker Studio)
\item Plan d'action pour le trimestre suivant
\end{enumerate}

\textbf{Soutenance :}
\begin{itemize}
\item Présentation : 15 min
\item Démonstration des résultats (avant/après) : 5 min
\item Questions : 10 min
\end{itemize}

\textbf{Barème global :}
\begin{itemize}
\item Rapport écrit : /40
\item Présentation orale : /30
\item Résultats mesurables : /20
\item Gestion de projet : /10
\item \textbf{Total : /100}
\end{itemize}

%%%%%%%%%%%%%%%%%%%%%%%%%%%%%%%%%%%%%%%%%%%%%%%%%%%
\chapter{SEO Spécialisé : Local, E-commerce, Migration et A/B Testing}
%%%%%%%%%%%%%%%%%%%%%%%%%%%%%%%%%%%%%%%%%%%%%%%%%%%

\section{Objectifs}
\begin{itemize}
\item Configurer et optimiser une fiche Google Business Profile
\item Mettre en œuvre une migration SEO complète
\item Optimiser un site e-commerce pour le SEO
\item Réaliser des tests A/B SEO et mesurer l'impact
\end{itemize}

\section{Outils nécessaires}
\begin{itemize}
\item Google Business Profile (compte Google)
\item Google Search Console
\item Screaming Frog
\item Google Optimize (ou équivalent A/B testing)
\item Outil de suivi de position (Sheets ou rank tracker)
\end{itemize}

\section{Exercice 1 --- Configuration Google Business Profile (45 min)}

\textbf{Consigne :} Créez une fiche Google Business Profile complète pour un commerce local.

\subsection{Étape 1 --- Création de la fiche}
Créez (ou simulez) une fiche GBP avec les informations suivantes :

\textbf{Entreprise fictive :}
\begin{itemize}
\item Nom : Librairie Le Savoir
\item Adresse : 25 Rue des Arts, 75001 Paris
\item Téléphone : 01 42 00 00 00
\item Site web : \url{https://librairie-le-savoir.fr}
\item Catégorie principale : Librairie
\item Catégories secondaires : Café littéraire, Salon de lecture
\item Horaires : Lun-Sam 9h-19h, Fermé dimanche
\end{itemize}

\textbf{À compléter dans la fiche :}
\begin{itemize}
\item[$\square$] Description de l'entreprise (250 caractères max, avec mots-clés locaux)
\item[$\square$] Horaires spéciaux (jours fériés)
\item[$\square$] Attributs (WiFi gratuit, accessible PMR, livraison)
\item[$\square$] Photos (logo, couverture, intérieur, extérieur, produits)
\item[$\square$] Posts Google (3 posts : actualité, promotion, événement)
\item[$\square$] Questions/Réponses (5 questions fréquentes avec réponses)
\item[$\square$] Produits / Services (5 catégories)
\item[$\square$] Attributs de confirmation (masque, distanciation\ldots)
\end{itemize}

\subsection{Étape 2 --- Optimisation locale}
\begin{enumerate}
\item Quels mots-clés locaux cibler dans la description ? (5 mots-clés)
\item Comment choisir les catégories secondaires ?
\item Quelle est la stratégie de photos (nombre, type, fréquence) ?
\item Comment gérer les attributs ``Confirmé par le propriétaire'' ?
\end{enumerate}

\subsection{Étape 3 --- Stratégie d'avis}
\begin{enumerate}
\item Rédigez 3 réponses à des avis (1 positif, 1 neutre, 1 négatif)
\item Proposez un processus pour obtenir 10 nouveaux avis par mois
\item Comment répondre à un avis négatif sans se mettre en danger ?
\item Quel est l'impact des réponses aux avis sur le classement local ?
\end{enumerate}

\section{Exercice 2 --- Migration SEO : Simulation complète (60 min)}

\textbf{Consigne :} Vous migrez le site \url{https://ancien-site.com} vers \url{https://nouveau-site.com}.

\subsection{Phase 1 --- Cartographie des URLs (15 min)}

\textbf{Anciennes URLs $\rightarrow$ Nouvelles URLs :}

\begin{longtable}{|l|l|c|}
\hline
\textbf{Ancienne URL} & \textbf{Nouvelle URL} & \textbf{Type (301/410)} \\
\hline
\verb|/index.php| & \verb|/| & \\
\hline
\verb|/article.php?id=1| & \verb|/blog/seo-guide| & \\
\hline
\verb|/produit.php?cat=5&id=12| & \verb|/produits/chaussures| & \\
\hline
\verb|/about-us.html| & \verb|/a-propos| & \\
\hline
\verb|/old-category/| & \verb|/categorie/| & \\
\hline
\verb|/blog/ancien-article-importants.html| & \verb|/blog/nouvel-article| & \\
\hline
\verb|/page-supprimee.html| & (supprimée) & \\
\hline
\verb|/tag/| & (supprimée) & \\
\hline
\end{longtable}

\begin{enumerate}
\item Complétez le tableau avec les types de redirection appropriés
\item Quels sont les risques d'une redirection 302 au lieu de 301 ?
\end{enumerate}

\subsection{Phase 2 --- Implémentation technique (20 min)}

\textbf{Règles de réécriture (Apache .htaccess) :}
\begin{lstlisting}
# Regles de redirection 301
RewriteEngine On
RewriteRule ^index\.php$ https://nouveau-site.com/ [L,R=301]
RewriteRule ^article\.php\?id=1$ https://nouveau-site.com/blog/seo-guide [L,R=301]
# Completez les regles manquantes...

# Gestion des erreurs 404 personnalisees
ErrorDocument 404 /404.html
\end{lstlisting}

\textbf{Questions :}
\begin{enumerate}
\item Comment gérer les paramètres d'URL (?session=, ?utm\_source=) ?
\item Faut-il conserver l'ancien robots.txt ? Le modifier ?
\item Comment mettre à jour le sitemap XML pendant la migration ?
\end{enumerate}

\subsection{Phase 3 --- Suivi post-migration (15 min)}

\textbf{Checklist quotidienne (J+1 à J+30) :}

\begin{longtable}{|l|l|l|}
\hline
\textbf{Jour} & \textbf{Action} & \textbf{Outil} \\
\hline
J+1 & & \\
\hline
J+2 & & \\
\hline
J+3 & & \\
\hline
J+7 & & \\
\hline
J+14 & & \\
\hline
J+30 & & \\
\hline
\end{longtable}

\textbf{KPIs de suivi :}
\begin{itemize}
\item Trafic organique : \_\_\_\_\_ vs \_\_\_\_\_ (avant/après)
\item Pages indexées : \_\_\_\_\_ vs \_\_\_\_\_
\item Taux d'erreur 4xx : \_\_\_\_\_ vs \_\_\_\_\_
\item Positions mots-clés principaux : \_\_\_\_\_ vs \_\_\_\_\_
\end{itemize}

\subsection{Phase 4 --- Gestion de crise (10 min)}

\textbf{Consigne :} Vous recevez une alerte : 2 jours après la migration, le trafic a chuté de 60\%.

\begin{enumerate}
\item Quelles sont les 3 premières choses que vous vérifiez ?
\item Quels outils utilisez-vous pour chaque vérification ?
\item Quelles décisions prenez-vous à J+2 ?
\item À quel moment décidez-vous de ``revenir en arrière'' ?
\end{enumerate}

\section{Exercice 3 --- Site E-commerce : Audit et Optimisation (60 min)}

\textbf{Consigne :} Vous auditez et optimisez un site e-commerce de 500 produits.

\subsection{Partie A --- Architecture du site}

\textbf{Structure actuelle :}
\begin{lstlisting}
Accueil
+-- Categories (12)
|   +-- Sous-categories (48)
|   |   +-- Pages produits (500)
|   |   +-- Pages variantes (2000)
+-- Blog (30 articles)
+-- Panier
+-- Compte
+-- Pages legales
\end{lstlisting}

\textbf{Questions :}
\begin{enumerate}
\item La profondeur de clics est-elle optimale ? (combien de clics pour un produit ?)
\item Proposez une architecture plus plate (flat architecture).
\item Comment gérer le contenu dupliqué entre les variantes (couleurs, tailles) ?
\item Faut-il indexer les pages de sous-catégories vides ?
\end{enumerate}

\subsection{Partie B --- Filtres et facettes}

Le site a des filtres : prix, couleur, taille, marque, matière.

\textbf{Problème :} Chaque filtre crée une URL unique $\rightarrow$ 15 000 URLs générées.

\begin{enumerate}
\item Quelles sont les bonnes pratiques pour gérer les URLs de filtres ?
\item Faut-il utiliser \texttt{noindex} ou \texttt{canonical} sur les pages de filtres ?
\item Qu'est-ce que le ``contenu fin'' (thin content) des pages de filtres et comment l'éviter ?
\item Proposez une stratégie de gestion des facettes SEO-friendly.
\end{enumerate}

\subsection{Partie C --- Page catégorie optimisée}

Créez une page catégorie modèle pour la catégorie ``Chaussures de running''.

\textbf{Éléments obligatoires :}
\begin{itemize}
\item[$\square$] Title tag optimisé (50-60 car.)
\item[$\square$] Meta description (150-160 car.)
\item[$\square$] H1 contenant le mot-clé principal
\item[$\square$] Texte de catégorie (150-200 mots en haut de page)
\item[$\square$] Grille produits avec balisage schema.org
\item[$\square$] Fil d'Ariane (BreadcrumbList JSON-LD)
\item[$\square$] Filtres avec gestion SEO
\item[$\square$] Pagination avec \texttt{rel="next"/"prev"} (ou alternative)
\item[$\square$] Liens vers sous-catégories
\item[$\square$] Images catégorie optimisées
\end{itemize}

\subsection{Partie D --- Pages produits en masse}

\begin{enumerate}
\item Comment générer des descriptions uniques pour 500 produits sans créer de contenu dupliqué ?
\item Quelle est la structure de données structurées (Product, Offer, AggregateRating) ?
\item Comment gérer les produits hors stock ? Les supprimer ? Les noindexer ?
\item Quel est l'impact des avis et du classement par note sur le SEO ?
\end{enumerate}

\section{Exercice 4 --- A/B Testing SEO (45 min)}

\textbf{Consigne :} Concevez et analysez des tests A/B pour améliorer le SEO.

\subsection{Test 1 --- Title Tag}

\textbf{Hypothèse :} Un title plus long avec des mots-clés additionnels améliore le CTR.

\begin{longtable}{|l|l|}
\hline
\textbf{Version} & \textbf{Title} \\
\hline
Contrôle (A) & \texttt{Chaussures de running Homme | RunTech} \\
\hline
Test (B) & \texttt{Chaussures de running Homme - RunTech Ultra-Grip Pro | Trail et Route} \\
\hline
\end{longtable}

\textbf{Protocole :}
\begin{enumerate}
\item Combien de pages inclure dans le test ? (justifiez)
\item Quelle est la durée minimale du test ?
\item Quels KPIs mesurer ?
\item À partir de quel seuil de significativité concluez-vous ?
\end{enumerate}

\subsection{Test 2 --- Meta Description}

\textbf{Hypothèse :} Une meta description avec une question + appel à l'action améliore le CTR.

\begin{longtable}{|l|l|}
\hline
\textbf{Version} & \textbf{Meta Description} \\
\hline
Contrôle (A) & \texttt{Découvrez nos chaussures de running Homme. Grande sélection de modèles RunTech. Livraison rapide.} \\
\hline
Test (B) & \texttt{Vous cherchez des chaussures de running ? Essayez l'Ultra-Grip Pro. Testée par 234 coureurs $\bigstar$ 4.7/5. Livraison offerte dès 99\texteuro{}.} \\
\hline
\end{longtable}

\textbf{Questions :}
\begin{enumerate}
\item Quels sont les risques du test A/B sur les meta descriptions ?
\item Google réécrit les meta descriptions. Comment savoir si votre test est annulé par Google ?
\item Comment segmenter les résultats par type de requête ?
\end{enumerate}

\subsection{Test 3 --- Structure de page}

\textbf{Hypothèse :} Ajouter un sommaire interactif (Table of Contents) en haut de l'article réduit le taux de rebond et améliore le temps passé.

\begin{enumerate}
\item Concevez le protocole de test
\item Quels sont les biais potentiels ?
\item Comment mesurer l'impact SEO (pas seulement UX) ?
\end{enumerate}

\subsection{Analyse de résultats A/B}

\begin{longtable}{|l|c|c|c|c|}
\hline
\textbf{Métrique} & \textbf{Contrôle (A)} & \textbf{Test (B)} & \textbf{Variation} & \textbf{Significatif ?} \\
\hline
Impressions & 12 500 & 13 200 & +5.6\% & p=0.12 \\
\hline
Clics & 375 & 468 & +24.8\% & p=0.03 \\
\hline
CTR & 3.0\% & 3.55\% & +18.3\% & p=0.02 \\
\hline
Position moyenne & 8.2 & 8.5 & $-0.3$ & p=0.45 \\
\hline
Taux de rebond & 62\% & 58\% & $-4$~pts & p=0.08 \\
\hline
Temps sur page & 2m15s & 2m45s & +22\% & p=0.04 \\
\hline
\end{longtable}

\textbf{Questions :}
\begin{enumerate}
\item Quelle est votre conclusion sur le test B ?
\item Quels résultats sont statistiquement significatifs (p $<$ 0.05) ?
\item La position moyenne baisse légèrement. Faut-il s'inquiéter ?
\item Implémenteriez-vous la version B sur l'ensemble du site ? Pourquoi ?
\item Quelles sont les limites de ce test ?
\end{enumerate}

\section{Exercice 5 --- SEO pour CMS : Configuration WordPress (30 min)}

\textbf{Consigne :} Configurez WordPress pour un SEO optimal.

\textbf{Installation neuve de WordPress. À configurer :}

\subsection{Partie A --- Réglages WordPress}
\begin{itemize}
\item[$\square$] Permaliens : structure personnalisée \texttt{/\%postname\%/}
\item[$\square$] Titre du site : optimisé SEO
\item[$\square$] Slogan : pertinent ou masqué
\item[$\square$] Visibilité : indexation autorisée
\item[$\square$] Lecture : page d'accueil statique (page ``Accueil'')
\item[$\square$] Discussion : fermer les commentaires après 30 jours
\item[$\square$] Médias : organiser par dossiers (mois)
\end{itemize}

\subsection{Partie B --- Plugins SEO}

\begin{longtable}{|l|l|l|}
\hline
\textbf{Plugin} & \textbf{Fonction} & \textbf{Configuration clé} \\
\hline
Yoast SEO / Rank Math & & \\
\hline
WP Rocket / W3 Total Cache & & \\
\hline
Smush / ShortPixel & & \\
\hline
Redirection & & \\
\hline
Wordfence / Sucuri & & \\
\hline
\end{longtable}

\begin{enumerate}
\item Complétez le tableau
\item Quels plugins éviter (bloat, conflits, sécurité) ?
\item Y a-t-il des plugins à ne JAMAIS installer ? Lesquels ?
\end{enumerate}

\subsection{Partie C --- Configuration avancée}
\begin{enumerate}
\item Générez et soumettez le sitemap XML
\item Configurez les données structurées (Organization, Article, BreadcrumbList)
\item Optimisez la base de données (requêtes, index, cache)
\item Mettez en place un CDN (Cloudflare, etc.)
\item Configurez la pagination (rel next/prev ou équivalent)
\end{enumerate}

\section{Exercice 6 --- Rapport d'audit SEO Spécialisé (30 min)}

\textbf{Consigne :} Rédigez un rapport d'audit pour un site e-commerce local.

\textbf{Données du site :}
\begin{itemize}
\item Site : \url{quincaillerie-en-ligne.fr}
\item Type : E-commerce de quincaillerie (3 000 produits)
\item Zone : France métropolitaine + Belgique
\item Trafic mensuel : 18 000 visites
\item Taux de conversion : 1.8\%
\item Panier moyen : 65 \texteuro{}
\end{itemize}

\textbf{Structure du rapport :}

\begin{enumerate}
\item \textbf{SEO Local} (10 points)
   \begin{itemize}
   \item Fiche Google Business Profile
   \item Citations et cohérence NAP
   \item Avis clients
   \item Pages locales (magasins)
   \end{itemize}

\item \textbf{SEO E-commerce} (10 points)
   \begin{itemize}
   \item Architecture du site
   \item Pages catégories
   \item Fiches produits
   \item Gestion des variantes et filtres
   \end{itemize}

\item \textbf{SEO Technique} (10 points)
   \begin{itemize}
   \item Vitesse et Core Web Vitals
   \item Mobile-friendliness
   \item Données structurées
   \item Internationalisation (France/Belgique)
   \end{itemize}

\item \textbf{Plan d'action priorisé} (10 actions)
\end{enumerate}

\section{Rendu final}

\textbf{À rendre :}
\begin{enumerate}
\item Fiche Google Business Profile complète (PDF ou capture d'écran)
\item Plan de migration SEO (avec mapping URLs + règles .htaccess)
\item Audit e-commerce complet (checklist de 30 points)
\item Protocole et résultats de test A/B
\item Configuration WordPress optimisée
\end{enumerate}

\textbf{Barème :}
\begin{itemize}
\item Google Business Profile : /10
\item Plan de migration : /10
\item Audit e-commerce : /10
\item A/B Testing : /10
\item Configuration CMS : /10
\item Total : /50
\end{itemize}

\end{document}
